\documentclass[exercice,a4]{thedocument}%
\startdatakeys
    \source{24GENMATMEAG3}
    \BO{2026}
    \cycle{4}
    \competencesDuSocle{CA,RA}
    \objectifsApprentissage{C4A1aa,C4A2ad,C4B2cb,C4D2aa,C4D3ac}
    \motsCles{facteurs premiers,probabilité,équation,volume,vitesse}
\enddatakeys
\begin{document}%

\medskip

Pour chacune des affirmations suivantes, dire si elle est vraie ou fausse.

Toutes les réponses devront être justifiées.

\medskip

\begin{enumerate}
\item \textbf{Affirmation 1}

La décomposition en produit de facteurs premiers du nombre $260$ est $4 \times 5 \times 13$.
\item \textbf{Affirmation 2}

Une urne opaque contient des boules indiscernables au toucher: 3 boules blanches, 4 boules jaunes et 8 boules rouges.

On pioche au hasard une boule dans cette urne et on note sa couleur.

Une autre urne opaque contient des boules indiscernables au toucher : 1 boule marquée de la lettre A, 1 boule marquée de la lettre B et 3 boules marquées de la lettre C.

On pioche au hasard une boule dans cette urne et on note la lettre obtenue.

La probabilité d'obtenir une boule de couleur rouge est supérieure à la probabilité d'obtenir une boule marquée de la lettre C.
\item \textbf{Affirmation 3}

La solution de l'équation $7x + 5 = 2x - 2$ est $-1,4$.
\item \textbf{Affirmation 4}

On empile 10 pièces cylindriques de $1,9$~cm de diamètre et de $0,2$~cm de hauteur. Le volume du cylindre, arrondi à l'unité, formé par les $10$~pièces est de $6$~cm$^3$.

\emph{Rappel} : le volume d'un cylindre de rayon $R$ et de hauteur $h$ est égal à $\pi \times R^2 \times h$.
\item \textbf{Affirmation 5}

Un éléphant qui court à une vitesse de 5 m/s est plus rapide qu'un cochon qui se déplace à une vitesse de 17 km/h.
\end{enumerate}

\newpage
\showthesource
\end{document}
